%==============================================================================
% writing_sample.tex -- revised draft
%
% Red text (\todo{...}) marks what still needs writing, or numbers to take
% from the tables and the log produced by master.do. The final version should
% contain no \todo at all: search for it before submitting.
%
% To compile in Overleaf: upload this file, create folders tables/ and
% figures/, and upload into them the contents of output/tables and
% output/figures from the repository.
%==============================================================================
\documentclass[11pt]{article}
\usepackage[margin=1in]{geometry}
\usepackage{amsmath,amssymb}
\usepackage{booktabs}
\usepackage{graphicx}
\usepackage{setspace}
\usepackage{xcolor}
\usepackage[round]{natbib}
\usepackage[hidelinks]{hyperref}
\onehalfspacing

\newcommand{\todo}[1]{\textcolor{red}{[#1]}}

\title{Pension Receipt and Depression among Older Adults in India\thanks{This
paper began as a term paper for Econometrics I at the Indian Statistical
Institute, Delhi, taught by Abhiroop Mukhopadhyay, and has since been revised.
Replication code: \todo{GitHub link}. \todo{Acknowledgments, if any.} All
errors are my own.}}
\author{Aishwarya Vardhan Soni\\ Indian Statistical Institute, Delhi}
\date{\todo{Month} 2026}

\begin{document}
\maketitle

\begin{abstract}
\noindent\todo{Write 120--150 words, in this order. (1) Question and data:
LASI Wave 1 (2017--18), 29,627 adults aged 60 and above. (2) Outcome: probable
major depression from the CIDI-SF. (3) The preferred estimate (Table 2,
column 5) with its 95\% confidence interval, against a weighted prevalence of
9.2\%, and what it rules out. (4) The selection finding from Tables 1 and 3:
how recipients of work-related and social pensions differ, and what the
separate estimates show. (5) One sentence on what this implies for reading
cross-sectional evidence on pensions and mental health.}
\end{abstract}

\medskip
\noindent\textbf{Keywords:} pensions, depression, ageing, social protection, India\\
\noindent\textbf{JEL codes:} H55, I12, I38, J14

\newpage

%------------------------------------------------------------------------------
\section{Introduction}
%------------------------------------------------------------------------------

As life expectancy in India rises and the share of older adults grows, income
security, health, and well-being in later life have become central questions
for public policy. Pensions are the main instrument designed to support older
citizens: they provide a stable income after work ends and reduce
vulnerability in old age. Beyond this financial role, pension income may also
affect mental health, by lowering financial stress and reducing dependence on
family members. \todo{Add one concrete fact: the share of adults aged 60+ who
receive any pension, and how much of that is social rather than work-related
pensions (Table~\ref{tab:desc}).}

Depression among older adults is common but receives little policy attention.
In later life it is linked to declining physical health, social isolation,
widowhood, reduced mobility, and economic insecurity. In LASI, 9.2 percent of
adults aged 60 and above meet the criteria for probable major depression. For
many older Indians, especially in poorer households, mental well-being depends
not only on health but on household structure, social status, and whether they
have a regular income, which makes pension receipt a natural candidate for
study.

\todo{Paragraph 3: the question and why it is hard. Pensions could reduce
depression, but receipt is not random. Work-related pensions go to former
formal-sector workers, who are more educated, more urban, and better off.
Social pensions are targeted: the national old-age pension requires a
below-poverty-line household, and the widow pension requires widowhood, both of
which predict depression. Selection therefore runs in opposite directions for
the two kinds of pension, and a pooled comparison mixes the two biases.}

\todo{Paragraph 4: what the paper does and finds, with numbers. Data and
outcome in one sentence; the pooled estimate and its confidence interval; what
it rules out relative to prevalence; the separate estimates for work-related
and social pensions (Table~\ref{tab:type}); what happens when ration-card type,
the eligibility rule, is controlled for; robustness in one clause.}

\todo{Paragraph 5: contribution. Relate the paper to \citet{CaseDeaton1998} on
South Africa's old-age pension; to \citet{Kaushal2014} and
\citet{UnnikrishnanImai2020} on India's old-age pension and household welfare;
and to \citet{GalianiGertlerBando2016}, who use age and geographic eligibility
cutoffs in Mexico and find that a non-contributory pension lowers depression
scores. Say in one or two sentences what this paper adds: evidence from India's
national ageing survey on how selection into pensions shapes naive
comparisons. End with a one-sentence roadmap.}

%------------------------------------------------------------------------------
\section{Data}
%------------------------------------------------------------------------------

\subsection{The Longitudinal Ageing Study in India}

The analysis uses Wave 1 of the Longitudinal Ageing Study in India (LASI),
fielded in 2017--18 \citep{IIPS2020report, IIPS2020guide}. LASI is a nationally
representative survey of adults aged 45 and above and their spouses, covering
every state and union territory except Sikkim. Its questionnaire is harmonised
with the Health and Retirement Study and covers demographics, work and
retirement, pensions and social welfare schemes, health, and mental health.
Because it records depressive symptoms and pension receipt for the same
respondents, it is well suited to this question.

\subsection{Sample}

The population of interest is adults aged 60 and above, the group to whom
LASI's social-pension questions are addressed. Wave 1 interviewed 72,250
individuals, of whom 31,464 were aged 60 or above. I exclude 818 respondents who
were not administered the depression module (both screening questions missing,
most likely proxy interviews), 1,018 with missing caste, and one with missing
religion. The estimation sample contains 29,627 respondents in 22,700
households.

\subsection{Measuring depression}

The outcome is an indicator of probable major depression constructed from the
Composite International Diagnostic Interview -- Short Form
\citep[CIDI-SF;][]{Kessler1998}, administered in questions MH201--MH223. The
module has two screening pathways. The first begins with MH201, which asks
whether the respondent felt sad, blue, or depressed for two weeks or more in a
row during the past year; the second begins with MH214, which asks about a
two-week period of lost interest in most things. In each pathway, the symptom
questions are asked only if the episode lasted most of the day, nearly every
day (MH202--MH203 and MH215--MH216). The depressed-mood pathway records seven
symptoms (MH204--MH211) and the loss-of-interest pathway six (MH217--MH222);
Appendix~\ref{app:construction} lists them.

A respondent is classified as depressed if either pathway records three or
more symptoms:
\begin{equation*}
  \text{Depressed}_i = \mathbf{1}\big[\max(\text{MoodSymptoms}_i,\ \text{InterestSymptoms}_i) \geq 3\big].
\end{equation*}
Because the symptom questions are reached only through a positive screen that
meets the duration criteria, this preserves the two-stage structure of the
CIDI-SF and avoids classifying respondents as depressed on the basis of
isolated symptoms such as fatigue or poor sleep. Respondents who were never
administered the module are treated as missing rather than as not depressed.
In the estimation sample, 7.6 percent of respondents are classified as
depressed, or 9.2 percent using LASI's national weights.

\subsection{Measuring pension receipt}

LASI records pensions in two places. The work, retirement and pension module
asks respondents who work in, or have retired from, the organized sector
whether they currently receive a pension or expect one in future (WE412). The
social welfare module asks every respondent aged 60 and above, scheme by
scheme, whether they are aware of it (SW201) and whether they currently avail
its benefits (SW202). The cash schemes are the National Old Age Pension
Scheme, which requires a below-poverty-line household; the widow pension,
asked only of women; and other state or central pension schemes. I exclude the
Annapurna scheme, which provides food grains rather than cash.

The main treatment is current receipt of any pension:
\begin{equation*}
  \text{Pension}_i = \mathbf{1}\big[\text{WE412}_i = 1 \ \text{ or } \ \text{SW202}_{s,i} = 1 \text{ for some cash pension scheme } s\big].
\end{equation*}
Respondents who were not asked a question, because they never worked in the
organized sector or were unaware of a scheme, are coded as not receiving that
pension. A broader definition that adds respondents who expect a work-related
pension in future (WE412 $= 2$) is used only as a robustness check, since they
receive no pension income yet. Because the two kinds of pension reach very
different people (Section~\ref{sec:who}), I also report work-related and
social pensions separately, and social pensions by scheme.

\subsection{Controls}

The regressions control for characteristics that plausibly affect both pension
receipt and mental health: age and its square; sex; household structure;
urban residence; years of schooling (zero for respondents who never attended
school); caste (Scheduled Caste, Scheduled Tribe, Other Backward Class, and
none of these, with OBC as the reference group); religion (Hindu, Muslim,
Christian, Sikh, and other, grouping categories with very few respondents);
quintiles of monthly per capita consumption expenditure (MPCE); and state
fixed effects.

Household structure combines marital status and living arrangements into seven
categories: married and living with a spouse only, or with a spouse and
children; married with the spouse not co-resident; widowed and living alone,
with children, or with others; and other unmarried respondents. The two
variables cannot usefully enter separately, because no widowed respondent
lives with a spouse: with both in the regression, the contrast between widowed
and married respondents would be identified only from the 212 married
respondents who live apart from their spouse.

Finally, to capture the eligibility rule for the old-age pension directly, I
merge from the household questionnaire whether the household holds a
below-poverty-line (BPL) or Antyodaya ration card (CO301--CO302).

%------------------------------------------------------------------------------
\section{Empirical approach}\label{sec:method}
%------------------------------------------------------------------------------

I estimate linear probability models of the form
\begin{equation}
  100 \times \text{Depressed}_i = \beta\,\text{Pension}_i + \mathbf{X}_i'\gamma + \mu_{s(i)} + \varepsilon_i ,
\end{equation}
where $\mathbf{X}_i$ contains the controls described above and $\mu_{s(i)}$ is
a fixed effect for respondent $i$'s state. Scaling the outcome by 100 makes
$\beta$ the difference in the probability of depression, in percentage points,
between pension recipients and otherwise similar non-recipients.
Table~\ref{tab:main} adds the controls in four steps, from no controls to the
preferred specification in column~5. State fixed effects absorb all
state-level differences, including in pension coverage, benefit generosity,
and health systems, so the preferred estimate compares recipients and
non-recipients within the same state.

Three choices deserve comment. First, MPCE and household structure are
measured after pension receipt and may partly respond to it: a pension raises
household spending \citep{Kaushal2014} and can change whom an older person
lives with. Controlling for them may therefore absorb part of any effect
\citep{AngristPischke2009}, so Appendix Table~\ref{tab:robust} also reports
estimates without them. Second, several respondents are interviewed in each
household, so standard errors are clustered by household; clustering instead
by the village or urban ward from which LASI sampled households
\todo{gives similar standard errors --- confirm from Table~\ref{tab:robust},
column 3}. Because the outcome is binary, the errors are heteroskedastic by
construction, and the clustered standard errors are robust to this. Third, the
regressions are unweighted, since the controls include LASI's main design
variables (state, urban residence, and age); estimates using LASI's national
weights are reported as a robustness check \citep{SolonHaiderWooldridge2015}.
Descriptive statistics are weighted.

%------------------------------------------------------------------------------
\section{Results}
%------------------------------------------------------------------------------

\subsection{Who receives a pension?}\label{sec:who}

\input{tables/table1_descriptives.tex}

\todo{Two paragraphs on Table~\ref{tab:desc}. (i) Coverage: the share
receiving any pension, and how much of it is social pensions. (ii) Selection,
with numbers from the table: compare work-pension recipients with
non-recipients on schooling, urban residence, widowhood, BPL cards, and MPCE;
then do the same for social-pension recipients, and note whether the
differences run the other way, as the eligibility rules imply. (iii) Raw
depression prevalence in each column. Close with one sentence on what this
means for a pooled comparison.}

\subsection{Pension receipt and depression}

\input{tables/table2_main.tex}

\todo{Two paragraphs on Table~\ref{tab:main}. (i) The estimate in each column;
judge changes across columns against the standard errors rather than calling
them large or small. (ii) The preferred estimate (column 5) with its 95\%
confidence interval, from the \texttt{lincom} output at the end of the log,
expressed relative to prevalence: for example, ``I can rule out that pension
recipients are more than [x] percentage points less likely to be depressed,
about [y] percent of the sample mean.'' Do not interpret the coefficients on
the controls.}

\subsection{Work-related and social pensions}

\input{tables/table3_pension_type.tex}

\begin{figure}[htbp]
  \centering
  \includegraphics[width=0.8\textwidth]{figures/figure1_coefplot.pdf}
  \caption{Pension coefficients from the preferred specification. ``Any
  pension'' is from Table~\ref{tab:main}, column 5; the remaining estimates are
  from Table~\ref{tab:type}, column 2. Bars are 95\% confidence intervals.}
  \label{fig:coefplot}
\end{figure}

\todo{Two or three paragraphs on Table~\ref{tab:type} and
Figure~\ref{fig:coefplot}. (i) Work-related versus social pensions (column 1),
and the test of whether they differ (the \texttt{lincom} of their difference
in the log). (ii) By scheme (column 2): old-age versus widow pensions. (iii)
What happens when the BPL card is controlled for (column 3) and within BPL
households (column 4). (iv) Interpretation: say whether the pattern matches
what the selection in Section~\ref{sec:who} predicts, and be explicit that it
is consistent with selection rather than proof of it.}

\subsection{Robustness and heterogeneity}

\todo{One paragraph. The pension coefficient under survey weights,
village-level clustering, the broader definition including expected pensions,
without MPCE and household structure, and in a probit
(Appendix Table~\ref{tab:robust}). Then heterogeneity by sex, urban residence,
age 70+, and BPL status: cite the interaction terms and whether any is
significant (Appendix Table~\ref{tab:het}).}

%------------------------------------------------------------------------------
\section{Interpreting the estimates}
%------------------------------------------------------------------------------

These estimates are conditional associations, not causal effects. Pension
receipt is not randomly assigned: people qualify through their employment
histories or through programme eligibility rules, so recipients may differ from
non-recipients in ways the controls do not capture.

\todo{About half a page, organised as signed biases. (i) Work-related
pensions: lifetime formal employment brings higher income, financial literacy,
and social networks that MPCE and schooling capture imperfectly; this positive
selection pushes the coefficient toward lower depression among recipients.
(ii) Social pensions: eligibility requires poverty or widowhood, both
associated with depression, so negative selection pushes the coefficient
toward higher depression. (iii) Health: poor health can lead both to earlier
retirement, and hence pension receipt, and to depression. (iv) Measurement:
the CIDI-SF is a screening instrument subject to stigma and recall, and
misreported pension receipt attenuates the estimate. Close with two or three
sentences on designs that could identify a causal effect: eligibility
thresholds in age and BPL status, state variation in pension amounts as in
\citet{Kaushal2014}, and later LASI waves.}

%------------------------------------------------------------------------------
\section{Conclusion}
%------------------------------------------------------------------------------

\todo{One or two paragraphs: the question; the pooled estimate and what it
rules out; the opposite selection into the two kinds of pension and what the
separate estimates show; the implication that cross-sectional comparisons of
pension recipients cannot settle whether pensions improve mental health; and
one sentence on the most promising route to a causal answer.}

%------------------------------------------------------------------------------
\begin{thebibliography}{10}
%------------------------------------------------------------------------------

\bibitem[Angrist and Pischke(2009)]{AngristPischke2009}
Angrist, J.~D. and J.-S. Pischke (2009).
\newblock \emph{Mostly Harmless Econometrics: An Empiricist's Companion}.
\newblock Princeton, NJ: Princeton University Press.

\bibitem[Case and Deaton(1998)]{CaseDeaton1998}
Case, A. and A.~Deaton (1998).
\newblock Large cash transfers to the elderly in {South Africa}.
\newblock \emph{Economic Journal} 108(450), 1330--1361.

\bibitem[Galiani et~al.(2016)Galiani, Gertler, and Bando]{GalianiGertlerBando2016}
Galiani, S., P.~Gertler, and R.~Bando (2016).
\newblock Non-contributory pensions.
\newblock \emph{Labour Economics} 38, 47--58.

\bibitem[{IIPS}(2020a)]{IIPS2020guide}
International Institute for Population Sciences (2020a).
\newblock \emph{Data User Guide: Longitudinal Ageing Study in India (LASI)
Wave 1, 2017--18}.
\newblock Mumbai: IIPS.

\bibitem[{IIPS et~al.}(2020b)]{IIPS2020report}
International Institute for Population Sciences, NPHCE, MoHFW, Harvard T.~H.
Chan School of Public Health, and University of Southern California (2020b).
\newblock \emph{Longitudinal Ageing Study in India (LASI) Wave 1, 2017--18,
India Report}.
\newblock Mumbai: IIPS.

\bibitem[Kaushal(2014)]{Kaushal2014}
Kaushal, N. (2014).
\newblock How public pension affects elderly labor supply and well-being:
Evidence from {India}.
\newblock \emph{World Development} 56, 214--225.

\bibitem[Kessler et~al.(1998)Kessler, Andrews, Mroczek, Ustun, and
Wittchen]{Kessler1998}
Kessler, R.~C., G.~Andrews, D.~Mroczek, B.~Ustun, and H.-U. Wittchen (1998).
\newblock The {World Health Organization Composite International Diagnostic
Interview} Short-Form ({CIDI-SF}).
\newblock \emph{International Journal of Methods in Psychiatric Research}
7(4), 171--185.

\bibitem[Solon et~al.(2015)Solon, Haider, and Wooldridge]{SolonHaiderWooldridge2015}
Solon, G., S.~J. Haider, and J.~M. Wooldridge (2015).
\newblock What are we weighting for?
\newblock \emph{Journal of Human Resources} 50(2), 301--316.

\bibitem[Unnikrishnan and Imai(2020)]{UnnikrishnanImai2020}
Unnikrishnan, V. and K.~S. Imai (2020).
\newblock Does the old-age pension scheme improve household welfare?
{Evidence} from {India}.
\newblock \emph{World Development} 134.

\end{thebibliography}

%------------------------------------------------------------------------------
\appendix
%------------------------------------------------------------------------------

\section{Additional tables}

\input{tables/table4_robustness.tex}

\input{tables/table5_heterogeneity.tex}

\section{Variable construction}\label{app:construction}

\paragraph{Depression (CIDI-SF).} Depressed-mood pathway (screen MH201,
symptom items asked if the episode lasted most of the day, nearly every day,
MH202--MH203): (1) loss of interest in most things (MH204); (2) fatigue or low
energy (MH205); (3) loss or gain of appetite (MH206--MH207, counted once); (4)
trouble concentrating (MH208); (5) feeling worthless (MH209); (6) thoughts of
death (MH210); (7) trouble sleeping (MH211). Loss-of-interest pathway (screen
MH214, criteria MH215--MH216): (1) fatigue or low energy (MH217); (2) loss of
appetite (MH218); (3) trouble concentrating (MH219); (4) feeling worthless
(MH220); (5) thoughts of death (MH221); (6) trouble sleeping (MH222).

\paragraph{Pension receipt.} Work-related: WE412, ``Please tell me, whether
you are currently receiving pension or expect to receive pension in future,''
asked if the respondent currently works in or has retired from the organized
sector (WE004 $= 1$ or WE402 $= 1$); 1 = currently receiving, 2 = expected in
future, 3 = neither. Social: SW202, asked for each scheme the respondent is
aware of (SW201 $= 1$): (a) National Old Age Pension Scheme; (b) widow pension
(women only); (c) Annapurna (excluded, as it provides food grains); (d) and
(e) other pension schemes.

\paragraph{Ration card.} CO301 (household has a ration card) and CO302 (card
type: BPL, Antyodaya, APL, or other), from the household questionnaire. The
indicator equals one for BPL or Antyodaya cards and zero for APL, other, or no
card.

\end{document}
